\section{Algoritmo}

Estas notas estudian algoritmos. Un \concept{algoritmo} es una secuencia de pasos. Pero nos interesan los algoritmos que resuelven un problema de computación, entonces:

\begin{definicion}{Algoritmo}{algoritmo}
Secuencia de pasos que transforman un conjunto de valores de \concept{entrada} (\textit{input}) en un conjunto de valores de \concept{salida} (\textit{output}).
\end{definicion}

Un algoritmo es un objeto matemático, independientemente de que luego deba programarse con algún lenguaje en alguna máquina. Para estudiarlo como entidad matemática, un algoritmo se puede expresar en:
\begin{itemize}
  \item Lenguaje natural (español, inglés, etc.)
  \item Lenguaje formal (matemática o algún lenguaje de programación o lógico)
  \item Seudocódigo (una mezcla de los dos anteriores)
  \item Como programa de computador
  \item Como hardware (esto es difícil de imaginar en este punto, pero sí)
\end{itemize}

Para estudiar algoritmos, es frecuente escribirlos en seudocódigo.
